\documentclass[12pt]{article}
\usepackage[utf8]{inputenc}
\usepackage[T1]{fontenc}
\usepackage[margin=1in]{geometry}
\usepackage{setspace}
\usepackage{xcolor}
\usepackage[colorlinks=true, urlcolor=blue!60!black, linkcolor=blue!60!black]{hyperref}
\setlength{\parindent}{2em}
\setlength{\parskip}{1.25em}
\renewcommand{\baselinestretch}{1.0}
\begin{document}
\begin{center}
{\LARGE \textit{A.L.A. Schechter Poultry Corp. v. United States}, 295 U.S. 495 (1935)\par}
\vspace{0.8em}
{\large Case Brief --- Four-Page Edition\par}
\end{center}
\vspace{1em}

\section{Memory Jogger}

In the ``sick chicken case,'' a unanimous Court reversed a Brooklyn kosher slaughterhouse's convictions under a New Deal poultry code, holding that NIRA \S~3 unconstitutionally delegated legislative power to the President to adopt binding codes, including codes proposed by trade groups, and that the code's reach to purely local slaughtering and sales exceeded the commerce power (\href{https://supreme.justia.com/cases/federal/us/295/495/}{295 U.S. 495}).

\section{Detailed Case Facts}

The National Industrial Recovery Act of 1933 authorized the President, on application by ``truly representative'' trade associations, to approve ``codes of fair competition''---or prescribe them himself---if they would ``tend to effectuate the policy of this title.'' Violating an approved code in any transaction affecting interstate commerce was a misdemeanor punishable by up to \$500 per day. In 1934 President Roosevelt approved the Live Poultry Code for the New York metropolitan area, setting a 40-hour week, a 50-cent minimum wage, reporting duties, and trade-practice rules, chief among them ``straight killing,'' which barred retail buyers from selecting individual chickens.

The Schechter brothers ran wholesale kosher slaughterhouses in Brooklyn. Although 96\% of New York's live poultry came from other States, they bought birds in New York markets, slaughtered them in Brooklyn, and sold them within 24 hours only to local retailers; they sold nothing in interstate commerce.

\section{Procedural History}

Indicted on 60 counts in the Eastern District of New York, the Schechters were convicted on 19, including conspiracy, wage-and-hour, ``straight killing,'' and selling an unfit chicken (the source of the nickname). The Second Circuit reversed the wage-and-hour counts as beyond the commerce power but affirmed the rest. The Supreme Court granted cross-petitions and decided the case on May 27, 1935.

\section{Judicial Votes}

\textbf{9--0}: the remaining convictions reversed and the reversal of the wage-and-hour counts affirmed.
\begin{itemize}
\item \textbf{Opinion of the Court:} Chief Justice Hughes, joined by Justices Van Devanter, McReynolds, Brandeis, Sutherland, Butler, and Roberts.
\item \textbf{Concurrence:} Justice Cardozo, joined by Justice Stone.
\item \textbf{Dissents:} None.
\end{itemize}

\section{Holding}

NIRA \S~3 unconstitutionally delegates legislative power under Article I, \S~1, by authorizing binding codes of fair competition without intelligible standards, and the Live Poultry Code's application to the Schechters' intrastate slaughtering and sales exceeds the commerce power. The Court did not reach the due process claim.

\section{Analysis of Opinions}

\textbf{Hughes: threshold.} The Depression was grave, but ``[e]xtraordinary conditions do not create or enlarge constitutional power,'' and the codes were not voluntary; once approved they were criminal ``codes of laws'' binding all.

\textbf{Hughes: delegation.} Building on \textit{Panama Refining}, Hughes held that Congress may let executive officers make subordinate rules and find facts, but it must set the policy and standards. NIRA never defined ``fair competition,'' which matched neither the common-law tort nor the FTC Act's cabined standard with hearings, findings, and review; it swept in whatever trade groups proposed to serve the Act's broad and conflicting goals. Delegation to private associations was ``unknown to our law.'' The President's required finding was a mere ``statement of an opinion,'' and he could write codes himself, so \S~3 was ``without precedent,'' unlike upheld delegations to the ICC, the Radio Commission, or tariff officials.

\textbf{Hughes: commerce.} The poultry had come to ``a permanent rest within the State,'' so it was no longer in a stream of commerce, as in \textit{Swift} or \textit{Stafford}. Local wages, hours, and sales practices affected commerce only indirectly; the direct/indirect distinction was ``fundamental,'' for otherwise state authority would exist ``only by sufferance of the federal government.''

\textbf{Cardozo (with Stone).} The lone dissenter in \textit{Panama Refining} explained that there Congress gave discretion over the occasion for barring one defined act, but here it gave a ``roving commission to inquire into evils and, upon discovery, correct them.'' On commerce, ``[t]he law is not indifferent to considerations of degree''; otherwise ``there will be an end to our federal system.''

\textbf{Later developments.} \textit{Jones \& Laughlin}, \textit{Darby}, and \textit{Wickard} abandoned the direct/indirect test, though Hughes's warning resurfaced in \textit{Lopez}. With \textit{Panama Refining}, \textit{Schechter} remains the only decision striking a federal statute for delegating power to public officials; in \textit{FCC v. Consumers' Research} (2025), the Court reaffirmed the intelligible-principle test, with Justice Kagan calling \textit{Schechter} an ``object lesson'' and Justice Gorsuch's dissent relying on it.

\section{Examples: Future Applications}

\begin{enumerate}
\item \textbf{Same-Side (Invalid): Industry-Drafted AI Codes.} A statute letting the President approve criminally enforced AI codes drafted by leading labs, under conflicting goals and without APA procedures or review, tracks \textit{Schechter} and \textit{Carter Coal}.
\item \textbf{Same-Side (Invalid): Majority Producers Bind an Industry.} Letting two-thirds of producers and a majority of workers set binding wages and hours for every firm hands coercive power to private parties, as in \textit{Carter Coal}.
\item \textbf{Opposite-Side (Valid): Health-Based Air Standards.} Directing an agency to set air standards protecting public health with an adequate margin of safety, through notice-and-comment and judicial review, supplies an intelligible principle (\textit{Whitman}).
\item \textbf{Opposite-Side (Valid Commerce Power): Slaughterhouse Minimum Wage.} A federal minimum wage reaching a local slaughterhouse like the Schechters' is valid under \textit{Darby} and \textit{Wickard}'s aggregation principle.
\item \textbf{Fence-Sitter: Criminal ``Fair and Orderly Markets'' Rules.} An agency's criminally enforced rules to promote fair and orderly markets resemble upheld public-interest standards, but echo NIRA's open-ended goals that Justice Gorsuch has criticized.
\end{enumerate}

\section{Critique}

\textbf{Private delegation and manageability.} \textit{Schechter} reflected suspicion of corporatist self-regulation that let dominant trade-association members cartelize industries and crush small competitors like the Schechters. Justice Scalia observed in \textit{Mistretta} that deciding ``how much is too much'' is a question of degree ill suited to judicial enforcement, explaining why courts prefer statutory narrowing.

\textbf{Procedure as a proxy for standards.} Hughes contrasted NIRA with the FTC Act and ICC statutes: the absence of notice, hearings, records, and judicial review, later standardized by the Administrative Procedure Act, made NIRA's delegation uniquely dangerous.

\section{Key Quotations}

\begin{quote}
``The delegated power of legislation which has found expression in this code is not canalized within banks that keep it from overflowing. It is unconfined and vagrant \dots\ This is delegation running riot. No such plenitude of power is susceptible of transfer.''
\end{quote}
\noindent --- Justice Cardozo, concurring, 295 U.S. at 551, 553.

\begin{quote}
``Extraordinary conditions may call for extraordinary remedies. But the argument necessarily stops short of an attempt to justify action which lies outside the sphere of constitutional authority. Extraordinary conditions do not create or enlarge constitutional power.''
\end{quote}
\noindent --- Chief Justice Hughes, Opinion of the Court, 295 U.S. at 528.

\begin{quote}
``Congress cannot delegate legislative power to the President to exercise an unfettered discretion to make whatever laws he thinks may be needed or advisable for the rehabilitation and expansion of trade or industry.''
\end{quote}
\noindent --- Chief Justice Hughes, Opinion of the Court, 295 U.S. at 537--538.

\end{document}
